\documentclass[10pt,a4paper]{amsart}
\usepackage[margin=19mm]{geometry}
\usepackage{amsmath,amssymb,mathtools}
\usepackage[scr=rsfs,cal=euler]{mathalfa}
\usepackage{xfrac}
\usepackage[unicode,hidelinks,bookmarks=false]{hyperref}
\newcommand{\ie}{i.e.\ }
\newcommand{\cf}{cf.\ }
\newcommand{\resp}{respectively\ }
\newcommand{\Subsection}{\subsection}
\providecommand{\nobreakdash}{\nobreak\hbox{-}}
\setlength{\emergencystretch}{2em}
\multlinegap0pt
\usepackage{bm}
\usepackage{bbm}
\usepackage{dsfont}
\usepackage{xspace}
\usepackage{cancel}
\usepackage{mathtools}
\usepackage{amsthm}
\makeatletter
\@ifundefined{claim}{\newtheorem{claim}{Claim}}{}
\@ifundefined{lemma}{\newtheorem{lemma}[claim]{Lemma}}{}
\@ifundefined{thm}{\newtheorem{thm}[claim]{Theorem}}{}
\makeatother
\title[On some properties of free commutators with semicircular var]{On some properties of free commutators with semicircular variables}
\author{Mihai Popa, Kamil Szpojankowski}
\date{Source: arXiv:2511.13578v1 (CC BY 4.0)}
\begin{document}
\maketitle

\noindent\emph{Source and licence.} On some properties of free commutators with semicircular variables. Version arXiv:2511.13578v1 (CC BY 4.0), distributed under the Creative Commons Attribution 4.0 International licence, as recorded in the arXiv licence metadata for this exact version and shown on the abstract page https://arxiv.org/abs/2511.13578v1. Full rights notice: licenses/arxiv-2511.13578-CC-BY-4.0.txt.\par\smallskip

\noindent\textit{Editorial scope.} Counted unit: the Lemma whose source label is \texttt{lem:31} in the article (source0.tex lines 494--496, source label \texttt{lem:31}), delivered together with the article's own complete original proof of it (source0.tex lines 498--614, one \texttt{proof} environment of 117 lines). No mathematics is written, repaired or re-derived: statement and proof are the article's own text line for line. Required context retained uncounted: thm:product-entries (lines 470--485). Every definition, display and label of those lines is retained unchanged. Omitted from this package and not redistributed: 1--469, 486--493, 497--497, 615--2310. Nothing inside a retained excerpt is omitted or altered: every retained body is delivered exactly as the article's lines are written, so no omission receipt of any kind accompanies this package. Citations: the retained excerpts carry no citation command at all, so no bibliography entry of the article is transcribed and no reference is redistributed. Reference adaptation: none was needed and none was made; every reference inside the delivered unit resolves inside it, so no unresolved reference or citation is produced. Printed slips kept literal: no printed word, symbol, hypothesis or number of the counted statement or of its proof was altered. Layout and class: the article's own class is \texttt{amsart} with packages including hyperref, amssymb, amsmath, amscd, mathrsfs, yfonts, bbm, graphics, dsfont, footmisc, xcolor, tikz, float, caption; the wrapper is the lane's pinned \texttt{amsart} editorial wrapper at 10pt on A4, which restates the theorem-like environments the retained text needs and adds the article's own macro definitions that the retained text needs (none), with extra wrapper packages (bm, bbm, dsfont, xspace, cancel, mathtools). The delivered numbering is the unit's own. Assets: no retained span contains a figure environment, a graphics-inclusion command, a PDF-page inclusion, a table or a TikZ picture; no image file is copied into this package, the package declares no asset dependency and no image file is redistributed with it. Licence: version arXiv:2511.13578v1 of this article is distributed under the Creative Commons Attribution 4.0 International licence, as recorded in the arXiv licence metadata for this exact version and shown on the abstract page https://arxiv.org/abs/2511.13578v1. A full rights notice is delivered at licenses/arxiv-2511.13578-CC-BY-4.0.txt. Characters: every retained excerpt is pure ASCII, so the delivered engine, XeLaTeX with its default Latin Modern text font, reports no missing character.
\begin{thm}\label{thm:product-entries}
Consider the non-commutative probability space $(\mathcal{A}, \tau)$ and let
$(\kappa_\pi)_{\pi \in NC} $ be the correspondent free cumulants. Let $m $ be a positive integer and let $\sigma $ be an interval partition of $[n]$ with $m$ blocks, 
\[ \sigma = \big(1, \dots, i(1)\big), \big(i(1)+1, \dots, i(2)\big), 
\dots, 
\big(i(m-1)+1, \dots, n\big).
\]
For any variables $ a_1, a_2, \dots, a_n \in \mathcal{A}$  we have that
\[ 
\kappa_m \big( a_1\cdots a_{i(1)}, a_{i(1)+1} \cdots a_{i(2)}, \dots, a_{i(m-1)+1} \cdots a_n \big)
= 
\sum_{\substack{\pi\in NC(n)\\ \pi \vee \sigma = \mathbbm{1}_n}}
\kappa_{\pi}\big[a_1,a_2, \dots, a_n\big]
\]

\end{thm}

\begin{lemma}\label{lem:31}
 If $ s $ and $ x $ are free random variables such that $s $ is semicircular, then $ s $ and $ i[s, x] $ are \textbf{not} free.
\end{lemma}

\begin{proof}
	
	
It suffices to show that if 
 is even, then 
\begin{equation}\label{k:even}
\kappa_{4}\big(s, i[s, x], i[s, x], s \big) \neq 0. 
\end{equation} 

By linearity we have
\begin{align*}
	\kappa_{4}\big(s, i[s, x], i[s, x], s \big) =&-\kappa_4(s,sx,sx,s)+\kappa_4(s,sx,xs,s)\\ &+\kappa_4(s,xs,sx,s)-\kappa_4(s,xs,xs,s).
\end{align*} 
Observe that only second term in the sum above is non--zero. Indeed for $\kappa_4(s,sx,sx,s)$ consider which $s$ is paired with the last $s$, and observe that all possible pairing force the partition to violate the maximality property from Theorem \ref{thm:product-entries}. For $\kappa_4(s,sx,sx,s)$ observe that we can pair the first $s$ only with second or last $s$, otherwise we will separate second $s$ and $\kappa_1(s)=0$. But again joining first $s$ with second or last $s$ forces violation of maximality again. Similar argument applies to $\kappa_4(s,xs,xs,s)$.

For the second term similar analysis as above leads to
\begin{align*}
	\kappa_4(s,sx,xs,s)=\kappa_2(s)\kappa_2(x)\kappa_2(s)>0.
\end{align*} 
%
%and if 
%$ n $ 
%is odd, then
%\begin{equation}\label{k:odd} \kappa_{n+3}\big( s, i[s, x], i[s, x], s, i[s, x], \dots, i[s, x], s\big) \neq 0.
%\end{equation}
%
% Suppose first that $ n $ is even. Consider $ \sigma \in NC(2n+2) $ be given by \\
% $ \sigma = \{(1), (2, 3), (4, 5), \dots, (2n, 2n +1), (2n +2)\} $
% and assume that for all $ k \in [n] $ we have that
% $\{ a_{2k}, a_{2k+1} \} = \{ s, x\} $. Theorem \ref{thm:product-entries} gives that
% %
% \begin{align}
% \kappa_{n+2}\big( s, a_2a_3, a_4 a_5, \dots, a_{2n} a_{2n+1}, s \big) 
% =
% \sum_{\substack{\pi\in NC(2n+2)\\ \pi \vee \sigma = \mathbbm{1}_{2n+2}}}
% \kappa_{\pi}\big[s,a_2, \dots, a_{2n+1}, s\big]
% \end{align}
% Denote 
% $ \mathscr{X} = \{1, n+2\}\cup \{ k : a_k = s \}$ and 
% $ \mathscr{Y} =  \{ j: \ a_j = x \} $. 
% Since $ s $ and $ x $ are free, for
% \[
% \kappa_{\pi}\big[s,a_2, \dots, a_{2r+1}, s\big]
% \neq 0 
% \]
% the partition $\pi $ must not have blocks containing both $ s $ and $x $. Also, from the definition of 
% $ n $ and since $ s $ is semicircular, it follows that 
% $ \pi $ must have one block with $ n $ elements consisting on the set $ \mathscr{Y} $ and $ \frac{n+2}{2} $ blocks, each with $2$ elements from $\mathscr{X} $. Moreover, the blocks with elements from $\mathscr{X} $ have consecutive (mod. 2n+2) elements: if $(t, t+s )$ is such a block of $ \pi $, then
% the set
% $\mathscr{Y} \cap  \{ t+k;\  1 \leq k \leq s-1  \} $
%   forms at most one block of $\pi $ whose length is less than $ n$. 
% Therefore
% $\pi $
%  is either the non-crossing partition
%   $ 
%   \big\{ (4k-1, 4k):\  1 \leq k \leq \frac{n+1}{2}\big\}
%   \cup
%    \big\{(2, 5, 6, 9, 10,  \dots,  2n+1)  \big\} $
% or
% $\big\{(4k+1, 4k+2): \ 0 \leq k \leq \frac{n}{2} \big\} \cup \big\{ (3, 4, 7, 8, \dots, 2n)
% \big\}$,
%  But $ \pi \vee \sigma = \mathbbm{1}_{2n+2}$,
% so 
% $ \pi = \big\{(4k+1, 4k+2): \ 0 \leq k \leq \frac{n}{2} \big\} \cup \big\{ (3, 4, 7, 8, \dots, 2n)
% \big\} $
% which gives
% \begin{align*}
%  \kappa_{n+2}\big( s, i[s, x], \dots, i[s, x], s \big)
%  &=i^n 
%  \kappa_{n+2} \big( s, sx, -xs, sx, -xs, \dots,sx, -xs, s \big)
%  \\ 
%& = i^{n}(-1)^{\frac{n}{2}} 
% \kappa_\pi \big[ s, s, x, x, s, \dots, x, s, s \big]\\
% & = \kappa_n(x)\neq 0.
% \end{align*}
% 
% 
% If $ n $ is odd, consider 
% $ \sigma \in NC(2n +3) $
%  given by
%\[
%\sigma =\big\{ (1), (2, 3), (4, 5), (6), (7, 8), \dots, (2n+1, 2n+2), (2n+3)
%\big\},\]
%and, similarly to the setting above, asume that for all
%$ k \in \{ 2, 4\} \cup \{ 7, 9, \dots, 2n+1 \}$
%we have that 
%$\{ a_k, a_{k+1}\} = \{ s, Y\} $. Now Theorem \ref{thm:product-entries} gives that
%%
% \begin{align}\label{n:odd}
%\kappa_{n+3}\big( s, a_2a_3, a_4 a_5, &  s, a_7a_8,  \dots, a_{2n+1} a_{2n+2}, s \big) \\
%&=
%\sum_{\substack{\pi\in NC(2n+3)\\ \pi \vee \sigma = \mathbbm{1}_{2n+3}}}
%\kappa_{\pi}\big[s,a_2, \dots, a_5, s, a_7, \dots, a_{2n+2}, s\big]\nonumber
%\end{align}
%%
%Denoting 
%$ \mathscr{X} = \{ 1, 6, 2n+3\} \cup \{ k:\ a_k = s\}$
%and 
%$ \mathscr{Y} =\{ j:\ a_j = x \}$, 
%same argument as in the case $n $ even gives that the summand in the right hand side of (\ref{n:odd}) is zero unless $ \pi $ has one block with all elements of $ \mathscr{Y} $ and $\frac{n+3}{2} $ blocks each with 2 elements from $\mathscr{X} $ which must also be consecutive numbers mod. $2n+3 $. Together with the condition $ \pi \vee \sigma = \mathbbm{1}_{2n+3}$ 
%these give that 
%\[
% \pi = \{ (1, 2), (5, 6)\} \cup \{ (4k, 4k+1):\ 2 \leq k \leq \frac{n+1}{2}\} \cup \{ ( 3, 4, 7, 10, 11, \dots, 2n+1) \}
% \]
%which, furthermore gives that
%\begin{align*}
%\kappa_{n+3}\big( s, i[s, Y],& i[s, Y], s, i[s, Y], \dots, i[s, Y], s\big)\\
%= &i^n 
%\kappa_{n+3}
%\big(
%s, -Ys, sY, -Ys, s, sY, -Ys, \dots, sY, -Ys, s\big)\\
%=& i^n (-1)^{\frac{n+1}{2}}
% \kappa_\pi \big[s, Y, s, s, Y, s, s, Y, \dots, s, Y, Y, s, s\big]\\
% =&i^n (-1)^{\frac{n+1}{2}} \kappa_n(Y) \neq 0
%\end{align*}

\end{proof}

\end{document}
